\documentclass[../../main.tex]{subfiles}
\begin{document}

\begin{center}
  {\fontsize{48pt}{58pt}\selectfont \textbf{1962}\normalsize\textbf{年度}}
\end{center}

\begin{table}[htb]
  \centering
  \caption{出題テーマと難易度}
  \begin{tabular}{|c|p{9cm}|c|}
    \hline
    \textbf{問題番号} & \textbf{テーマ} & \textbf{難易度} \\
    \hline
    第1問 & 三角形の周上で中点を共有する線分の本数 & \\
    \hline
    第2問 & 定点を通る円群と中心の軌跡 & \\
    \hline
    第3問 & 円周上の動点と定点を通る直線 & \\
    \hline
    第4問 & 三次関数と二次関数の大小関係 & \\
    \hline
    第5問 & 三角関数の極値と極大点の配列 & \\
    \hline
    第6問 & 放物線の接線で囲まれる面積 & \\
    \hline
  \end{tabular}
\end{table}

\begin{table}[htb]
  \centering
  \caption{解答}
  \renewcommand{\arraystretch}{1.4}
  \begin{tabular}{|c|c|p{9cm}|}
    \hline
    \textbf{問題番号} & \textbf{小問} & \textbf{答え} \\
    \hline
    第1問 & -- & 各辺の中点を結ぶ三角形の外部では$1$本，辺上では$2$本，内部では$3$本（解答の図を参照） \\
    \hline
    第2問 & (1) & $(0,0)$と$\left(\dfrac{ab^2}{a^2+b^2},\dfrac{a^2b}{a^2+b^2}\right)$を通る \\
    \cline{2-3}
    & (2) & $2bx+2ay=ab$，ただし$x\ne0$ \\
    \hline
    第3問 & -- & $RS$は定点を通る（証明問題） \\
    \hline
    第4問 & -- & $c<a+b$のとき存在し，$c\ge a+b$のとき存在しない \\
    \hline
    第5問 & (1) & $0<a<1$ \\
    \cline{2-3}
    & (2) & 極大点は同一直線上に等間隔に並ぶ（証明問題） \\
    \hline
    第6問 & -- & $\dfrac{(a^2-4)^{3/2}-(a+4)(a-2)^2}{3}$ \\
    \hline
  \end{tabular}
\end{table}

\paragraph{試験形態}
6問．試験時間は要確認．

\end{document}
